\chapter{Introducción}\label{cap:introduccion}

\section{Pregunta y alcance}

La mortalidad por cáncer varía mucho entre los condados de Estados Unidos: en el fichero
analizado va de \res{min-valor} muertes por \num{100000} habitantes (\res{min-condado}) a
\res{max-valor} (\res{max-condado}). El objetivo del trabajo, fijado por la orientación del
proyecto, es construir un modelo de regresión lineal múltiple estimado por mínimos cuadrados
ordinarios que explique y prediga esa tasa (\var{TARGET\_deathRate}) a partir de las demás
características del condado, y acompañarlo de su diagnóstico completo.

Se distinguen dos usos del modelo, que requieren criterios distintos:
\begin{itemize}
  \item \textbf{Modelo de efectos} (\cref{cap:construccion,cap:modelo-final,cap:diagnostico,cap:sensibilidad}):
  estimar e interpretar la asociación de cada característica con la mortalidad, ajustada por
  las demás, con una inferencia válida (intervalos y contrastes con la cobertura y el nivel
  nominales).
  \item \textbf{Modelo predictivo} (\cref{cap:prediccion}): acertar en condados nuevos,
  evaluado con validación cruzada y una partición de prueba independiente.
\end{itemize}

Los datos son \emph{ecológicos}: cada observación es un condado, no una persona. Las
asociaciones estimadas describen condados y no deben trasladarse a individuos (falacia
ecológica, \cite{robinson1950}). Además, el diseño es observacional, por lo que los
coeficientes son asociaciones condicionales, no efectos causales.

\section{Correspondencia con los entregables}

La orientación pide seis entregables. La \cref{tab:entregables} indica dónde se responde cada uno.

\begin{table}[htbp]
\centering\small
\caption{Entregables de la orientación y su localización en el informe}\label{tab:entregables}
\begin{tabularx}{\linewidth}{c X l}
\toprule
& Entregable & Dónde \\
\midrule
(a) & Modelo final y pasos seguidos para construirlo & \cref{cap:construccion,cap:modelo-final} \\
(b) & Salida del software con $R^2$, $R^2$ ajustado y RMSE & \cref{sec:salida-spss} \\
(c) & Código & \cref{ap:reproducibilidad} \\
(d) & Diagnóstico: linealidad, independencia, homocedasticidad, normalidad, multicolinealidad & \cref{cap:diagnostico} \\
(e) & Interpretación del modelo & \cref{sec:interpretacion} \\
(f) & Atípicos, valores perdidos y variables categóricas & \cref{cap:datos,cap:ausentes} \\
\bottomrule
\end{tabularx}
\end{table}

\section{Cómo está hecho este informe}

Todo el análisis se ejecuta con un paquete de Python propio (\var{cancerstats}) a partir de
una configuración declarativa: cada decisión metodológica (umbrales, estimador, estrategia de
selección, interacciones candidatas\dots) es una opción con su justificación. Cada ejecución
queda registrada como una \emph{corrida} con su configuración, el entorno de software, la
huella del código y de los datos, el registro y todos los resultados. Las cifras, tablas y
figuras de este documento se generan automáticamente a partir de la última corrida
(\res{corrida-id}, del \res{corrida-fecha}); ninguna se ha copiado a mano. El
\cref{ap:reproducibilidad} describe cómo reproducirlo, la interfaz interactiva y la sintaxis
de SPSS equivalente.

La referencia teórica es \textcite{wasserman2004}, texto del curso, complementado con
\textcite{fox2016} y \textcite{wooldridge2010} para el diagnóstico y la inferencia robusta.
